% GENERATED FILE. Edit canonical lessons, then run npm run generate:books.
% canonical-source-hash: fnv1a64:f651c618eddaef81
% canonical-lessons: TE-C514-cevipotu, TE-C514-pantinoppi, TE-C514-vennunoppi, TE-C514-mulugu, TE-C514-murchapovu

\chapter{Aches}
\label{ch:te-aches}
\hlchaptermodality{514}

\begin{chapteropening}
\textbf{By the end of this chapter:} \emph{I can say where it hurts and that someone groaned or fainted.}

Complete the last lesson of chapter 514: I can say where it hurts and that someone groaned or fainted.
\end{chapteropening}

\section[\texorpdfstring{\te{చెవిపోటు} (cevipōṭu)}{చెవిపోటు (cevipōṭu)}]{\te{చెవిపోటు} (cevipōṭu) — an earache}

\label{lesson:TE-C514-cevipotu}



\begin{quote}
Before the new one: say the Telugu for paralysis, then the Telugu for asthma.
\end{quote}

\subsection*{You'll want to know: \te{చెవిపోటు}}
\textbf{\te{చెవిపోటు}} — \emph{cevipōṭu} — \textquotedblleft{}an earache\textquotedblright{}.

\begin{quote}
\te{చల్లని} \te{గాలి} \te{వల్ల} \te{పాపకు} \te{చెవిపోటు} \te{వచ్చింది}.
\end{quote}

\begin{quote}
\emph{callani gāli valla pāpaku cevipōṭu vaccindi.}
\end{quote}

\begin{quote}
The cold wind gave the little girl an earache.
\end{quote}

\te{చెవి}, ear, + \te{పోటు}, a throbbing pain.

\begin{cousinweb}
\textbf{\te{చెవిపోటు}} (\emph{cevipōṭu}) is native Dravidian: DEDR 1977 lists Telugu \emph{cevi}, \textquotedblleft{}ear\textquotedblright{}, with Tamil \emph{cevi}; \emph{pōṭu} is a stab of pain.
\end{cousinweb}

\begin{grammarlens}[title={What you've built}]
One more word for describing what you lived and felt.
\end{grammarlens}

\subsection*{Your turn}
\begin{itemize}
\raggedright
  \item \emph{Say it:} \emph{cevipōṭu}
  \item \emph{Say it:} \emph{cevipōṭu}, once more
  \item \emph{Say it:} \emph{callani gāli valla pāpaku cevipōṭu vaccindi.} — The cold wind gave the little girl an earache.
  \item \emph{Recall:} say the Telugu for paralysis, then the Telugu for asthma, then say \emph{cevipōṭu} again
\end{itemize}

\begin{tcolorbox}[breakable,colback=teal!4,colframe=teal!35!black,title={Before you move on}]
What does \te{చెవిపోటు} mean? (\textquotedblleft{}An earache\textquotedblright{}.) Say it once more.
\end{tcolorbox}
\section[\texorpdfstring{\te{పంటినొప్పి} (paṇṭinoppi)}{పంటినొప్పి (paṇṭinoppi)}]{\te{పంటినొప్పి} (paṇṭinoppi) — a toothache}

\label{lesson:TE-C514-pantinoppi}



\begin{quote}
Before the new one: say the Telugu for asthma, then the Telugu for an earache.
\end{quote}

\subsection*{You'll want to know: \te{పంటినొప్పి}}
\textbf{\te{పంటినొప్పి}} — \emph{paṇṭinoppi} — \textquotedblleft{}a toothache\textquotedblright{}.

\begin{quote}
\te{పంటినొప్పితో} \te{రాత్రంతా} \te{నిద్ర} \te{పట్టలేదు}.
\end{quote}

\begin{quote}
\emph{paṇṭinoppitō rātrantā nidra paṭṭalēdu.}
\end{quote}

\begin{quote}
With a toothache I could not sleep all night.
\end{quote}

\te{పన్ను}, tooth, takes the form \te{పంటి}- before another word.

\begin{cousinweb}
\textbf{\te{పంటినొప్పి}} (\emph{paṇṭinoppi}) is native Dravidian: \emph{pallu}, \emph{pannu}, \textquotedblleft{}tooth\textquotedblright{} (DEDR 3986), in its joining form \emph{paṇṭi-}, and \emph{noppi}, \textquotedblleft{}pain\textquotedblright{} (DEDR 3793).
\end{cousinweb}

\begin{grammarlens}[title={What you've built}]
One more word for describing what you lived and felt.
\end{grammarlens}

\subsection*{Your turn}
\begin{itemize}
\raggedright
  \item \emph{Say it:} \emph{paṇṭinoppi}
  \item \emph{Say it:} \emph{paṇṭinoppi}, once more
  \item \emph{Say it:} \emph{paṇṭinoppitō rātrantā nidra paṭṭalēdu.} — With a toothache I could not sleep all night.
  \item \emph{Recall:} say the Telugu for asthma, then the Telugu for an earache, then say \emph{paṇṭinoppi} again
\end{itemize}

\begin{tcolorbox}[breakable,colback=teal!4,colframe=teal!35!black,title={Before you move on}]
What does \te{పంటినొప్పి} mean? (\textquotedblleft{}A toothache\textquotedblright{}.) Say it once more.
\end{tcolorbox}
\section[\texorpdfstring{\te{వెన్నునొప్పి} (vennunoppi)}{వెన్నునొప్పి (vennunoppi)}]{\te{వెన్నునొప్పి} (vennunoppi) — a backache}

\label{lesson:TE-C514-vennunoppi}



\begin{quote}
Before the new one: say the Telugu for an earache, then the Telugu for a toothache.
\end{quote}

\subsection*{You'll want to know: \te{వెన్నునొప్పి}}
\textbf{\te{వెన్నునొప్పి}} — \emph{vennunoppi} — \textquotedblleft{}a backache\textquotedblright{}.

\begin{quote}
\te{రోజంతా} \te{కూర్చుని} \te{పనిచేయడం} \te{వల్ల} \te{నాకు} \te{వెన్నునొప్పి} \te{వచ్చింది}.
\end{quote}

\begin{quote}
\emph{rōjantā kūrcuni panicēyaḍaṁ valla nāku vennunoppi vaccindi.}
\end{quote}

\begin{quote}
Sitting and working all day gave me a backache.
\end{quote}

\te{వెన్ను}, the back or spine, + \te{నొప్పి}, pain.

\begin{cousinweb}
\textbf{\te{వెన్నునొప్పి}} (\emph{vennunoppi}) is native Dravidian: \emph{vennu}, \textquotedblleft{}the back\textquotedblright{} (DEDR 5488), and \emph{noppi}, \textquotedblleft{}pain\textquotedblright{} (DEDR 3793), with Tamil \emph{nō}, \textquotedblleft{}to ache\textquotedblright{}.
\end{cousinweb}

\begin{grammarlens}[title={What you've built}]
One more word for describing what you lived and felt.
\end{grammarlens}

\subsection*{Your turn}
\begin{itemize}
\raggedright
  \item \emph{Say it:} \emph{vennunoppi}
  \item \emph{Say it:} \emph{vennunoppi}, once more
  \item \emph{Say it:} \emph{rōjantā kūrcuni panicēyaḍaṁ valla nāku vennunoppi vaccindi.} — Sitting and working all day gave me a backache.
  \item \emph{Recall:} say the Telugu for an earache, then the Telugu for a toothache, then say \emph{vennunoppi} again
\end{itemize}

\begin{tcolorbox}[breakable,colback=teal!4,colframe=teal!35!black,title={Before you move on}]
What does \te{వెన్నునొప్పి} mean? (\textquotedblleft{}A backache\textquotedblright{}.) Say it once more.
\end{tcolorbox}
\section[\texorpdfstring{\te{మూలుగు} (mūlugu)}{మూలుగు (mūlugu)}]{\te{మూలుగు} (mūlugu) — to groan; to moan}

\label{lesson:TE-C514-mulugu}



\begin{quote}
Before the new one: say the Telugu for a toothache, then the Telugu for a backache.
\end{quote}

\subsection*{You'll want to know: \te{మూలుగు}}
\textbf{\te{మూలుగు}} — \emph{mūlugu} — \textquotedblleft{}to groan; to moan\textquotedblright{}.

\begin{quote}
\te{నొప్పితో} \te{రోగి} \te{రాత్రంతా} \te{మూలిగాడు}.
\end{quote}

\begin{quote}
\emph{noppitō rōgi rātrantā mūligāḍu.}
\end{quote}

\begin{quote}
The patient groaned all night with pain.
\end{quote}

Past (I …): \te{మూలిగాను}, \emph{mūligānu}.

The sound of pain or tiredness; also to grumble.

\begin{cousinweb}
\textbf{\te{మూలుగు}} (\emph{mūlugu}) is native Dravidian: DEDR 4896 lists Telugu \emph{mūl(u)gu}, \textquotedblleft{}to groan, to moan as from pain, to grumble\textquotedblright{}, with Kannada \emph{mūlugu} and Malayalam \emph{mūḷuka}, \textquotedblleft{}to groan\textquotedblright{}.
\end{cousinweb}

\begin{grammarlens}[title={What you've built}]
One more word for describing what you lived and felt.
\end{grammarlens}

\subsection*{Your turn}
\begin{itemize}
\raggedright
  \item \emph{Say it:} \emph{mūlugu}
  \item \emph{Say it:} \emph{mūlugu}, once more
  \item \emph{Say it:} \emph{noppitō rōgi rātrantā mūligāḍu.} — The patient groaned all night with pain.
  \item \emph{Recall:} say the Telugu for a toothache, then the Telugu for a backache, then say \emph{mūlugu} again
\end{itemize}

\begin{tcolorbox}[breakable,colback=teal!4,colframe=teal!35!black,title={Before you move on}]
What does \te{మూలుగు} mean? (\textquotedblleft{}To groan; to moan\textquotedblright{}.) Say it once more.
\end{tcolorbox}
\section[\texorpdfstring{\te{మూర్ఛపోవు} (mūrchapōvu)}{మూర్ఛపోవు (mūrchapōvu)}]{\te{మూర్ఛపోవు} (mūrchapōvu) — to faint}

\label{lesson:TE-C514-murchapovu}



\begin{quote}
Before the new one: say the Telugu for a backache, then the Telugu for to groan.
\end{quote}

\subsection*{You'll want to know: \te{మూర్ఛపోవు}}
\textbf{\te{మూర్ఛపోవు}} — \emph{mūrchapōvu} — \textquotedblleft{}to faint\textquotedblright{}.

\begin{quote}
\te{ఎండలో} \te{నిలబడి} \te{ఒక} \te{విద్యార్థి} \te{మూర్ఛపోయాడు}.
\end{quote}

\begin{quote}
\emph{eṇḍalō nilabaḍi oka vidyārthi mūrchapōyāḍu.}
\end{quote}

\begin{quote}
Standing in the sun, a student fainted.
\end{quote}

Past (I …): \te{మూర్ఛపోయాను}, \emph{mūrchapōyānu}.

\te{మూర్ఛ}, a faint or a fit, + \te{పోవు}, to go.

\begin{cousinweb}
\textbf{\te{మూర్ఛపోవు}} (\emph{mūrchapōvu}) is built from \te{మూర్ఛ}, from Sanskrit \emph{mūrchā}, \textquotedblleft{}fainting, stupor\textquotedblright{} (CDIAL 10242) and \emph{pōvu}, \textquotedblleft{}to go\textquotedblright{} (DEDR 4572).
\end{cousinweb}

\begin{grammarlens}[title={What you've built}]
One more word for describing what you lived and felt.
\end{grammarlens}

\subsection*{Your turn}
\begin{itemize}
\raggedright
  \item \emph{Say it:} \emph{mūrchapōvu}
  \item \emph{Say it:} \emph{mūrchapōvu}, once more
  \item \emph{Say it:} \emph{eṇḍalō nilabaḍi oka vidyārthi mūrchapōyāḍu.} — Standing in the sun, a student fainted.
  \item \emph{Recall:} say the Telugu for a backache, then the Telugu for to groan, then say \emph{mūrchapōvu} again
\end{itemize}

\begin{tcolorbox}[breakable,colback=teal!4,colframe=teal!35!black,title={Before you move on}]
What does \te{మూర్ఛపోవు} mean? (\textquotedblleft{}To faint\textquotedblright{}.) Say it once more.
\end{tcolorbox}
